\documentclass{standalone}
\usepackage{pstricks,pst-plot,pstricks-add}
\usepackage{pst-func}
\usepackage{pst-eucl}
\definecolor{cgrille}{gray}{0.3}
\usepackage{graphicx}
\usepackage{amsmath}
\usepackage{bm}
\newcommand{\degre}{\ensuremath{^\circ} }

%\uput{distance}[direction]{rotation}(x,y){texte}
%\pstMarkAngle[linecolor=red, arrows=->]{B}{A}{C}{$\gamma$} Permet de marquer les angles!
%	Option a mettre entre [] :
%		LabelSep : (1 unité par défaut) 
%		MarkAngleRadius : Rayon pour marquer les angles

\begin{document}
\psset{yunit=1cm,xunit=1cm
,plotpoints=1000,algebraic=true,
subgriddiv=1, griddots=5, gridlabels=0pt,
labels=all, labelFontSize=\tiny , ticks=all, ticksize=3pt , tickstyle=full,  Dy=1,Dx=1}
\begin{pspicture*}(0,0)(5.5,3)
%\psgrid[subgriddiv=1,griddots=7,gridlabels=0pt](0,0)(6,3)
%\psaxes[Dy=1, Dx=1, labelFontSize=\scriptstyle]{->}(0,0)(10,4)
\uput{0}[0]{0}(0,2.5){$\underline{1524}=\bm{8}\cdot 180+84$}
\psline{->}(0.5,1.8)(1.85,2.28)
\psline{->}(1.9,1.8)(2.7,2.28)
\uput{0}[0]{0}(0,1.5){$\phantom{1}\underline{180}=\bm{2}\cdot 84+\boxed{12}$}
\psline{->}(0.5,0.8)(1.75,1.25)
\psline{->}(2,0.75)(2.7,1.25)
\uput{0}[0]{0}(0,0.5){$\phantom{15}\underline{84}=\bm{7}\cdot \boxed{12}+0$}
\psline{->}(3,1.5)(3.5,1.5)
\uput{0}[0]{0}(3.6,1.5){\begin{minipage}{3cm}\footnotesize Dernier reste\\  non-nul\end{minipage}}
\end{pspicture*}
\end{document}
